\documentclass[12pt, letterpaper]{article}

% ── Packages ────────────────────────────────────────────────
\usepackage[margin=1in]{geometry}
\usepackage{setspace}
\usepackage{parskip}
\usepackage{titlesec}
\usepackage{graphicx}
\usepackage{amsmath, amssymb}
\usepackage{booktabs}
\usepackage{caption}
\usepackage{subcaption}
\usepackage{hyperref}
\usepackage{fancyhdr}
\usepackage{times}  % Times New Roman font
\usepackage{microtype}
\usepackage{enumitem}
\usepackage{float}  % [H] figure placement
\usepackage{multirow}

% ── Hyperlink Styling ────────────────────────────────────────
\hypersetup{
    colorlinks = true,
    linkcolor  = blue,
    citecolor  = blue,
    urlcolor   = blue,
}

% ── Section / Subsection Formatting ─────────────────────────
\titleformat{\section}{\large\bfseries}{}{0em}{}[\titlerule]
\titleformat{\subsection}{\normalsize\bfseries}{}{0em}{}
\titleformat{\subsubsection}{\normalsize\itshape}{}{0em}{}

% ── Header / Footer ─────────────────────────────────────────
\pagestyle{fancy}
\fancyhf{}
\rhead{\small CS 679 Pattern Recognition}
\lhead{\small Assignment 4}
\cfoot{\thepage}
\renewcommand{\headrulewidth}{0.4pt}

% ── Spacing ─────────────────────────────────────────────────
\onehalfspacing

% ============================================================
%  DOCUMENT
% ============================================================
\begin{document}

% ── COVER PAGE ──────────────────────────────────────────────
\begin{titlepage}
	\centering
	\vspace*{2cm}

	{\LARGE \textbf{CS 679 Pattern Recognition}}\\[0.5cm]

	{\Large \textbf{Assignment 4 - Gender Classification}}\\[0.5cm]

	\vfill

	{\large \textbf{Richard White}}\\[0.3cm]
	{\normalsize May 15, 2026}\\[1cm]

	\vfill

	% ── Academic Integrity Statement ────────────────────────
	\begin{minipage}{0.85\textwidth}
		\small
		\textit{``I declare that all material in this assignment is my own work except where there is clear
			acknowledgment or reference to the work of others. I understand that both my report and code may be
			subjected to plagiarism detection software, and fully accept all consequences if found responsible for
			plagiarism, as explained in the syllabus, and described in UNR’s Academic Standards Policy: UAM
			6,502.''}

		\vspace{1cm}
		\noindent\rule{6cm}{0.4pt}\\
		Richard White
	\end{minipage}

	\vspace{2cm}
\end{titlepage}

% ── TABLE OF CONTENTS ───
\tableofcontents
\newpage

% ============================================================
%  SECTION 1 -- THEORY
% ============================================================
\section{Theory}
\label{sec:theory}

% ============================================================
%  SECTION 2 -- RESULTS AND DISCUSSION
% ============================================================
\section{Results and Discussion}
\label{sec:results}

% Present and discuss your results. Include sample results, discuss major
% findings, and reference any figures or tables generated (i.e., using figure/table captions) to illustrate
% your results. For clarity, the results for each part must be presented in a separate subsection.

The experiments were conducted with a gender dataset containing 400 images from 400 people
with equal representation between males and females. Two image resolutions are used: 16x20
and 48x60 pixels. The dataset is split into a training, validation, and testing
sets and evaluated using a three-fold cross-validation.

\subsection{Experiments 1 \& 2}

In this set of experiments, we compare the performance of SVMs and Bayes classifier. The SVM is
implemented with a polynomial and RBF kernel. The Bayes classifier is implemented with a Gaussian
distribution and parameters are estimated using ML estimation on the training data and Mahalanobis
distance for classification. Equal priors $P(\omega_1) = P(\omega_2)$ are assumed in the
calculations since the dataset is balanced. In order to fairly compare between SVMs and Bayes,
the same eigen-cofficients are used between the two classifiers.

\paragraph{SVM.}
After performing the sweep, the best parameters for the polynomial kernel were found to be $C=0.1$ and $D=1$ for both resolutions except for the 48x60
resolution where $D=3$ performed best.
For the RBF kernel, the best parameters for each resolution were found to have one fold with $C=1$ and $\gamma=0.1$ and two folds with $C=10$ and $\gamma=0.1$.
The best parameters for each fold and resolution are shown in Tables \ref{tab:svm_bayes_16x20} and \ref{tab:svm_bayes_48x60}.
It's notable that RBF kernels generally performed better than polynomial kernels in the low resolution case, while in the high resolution case, the performance was more mixed.

\paragraph{Bayes Classifier.}
In comparison, the Bayes classifier performed significantly worse than the SVMs, with test errors ranging from 14.04\% to 18.80\% across folds and resolutions.
The Mahalanobis distance approach yielded higher error rates compared to the SVMs, indicating that the assumptions of Gaussian distributions may not hold well for
this dataset.

\paragraph{Why did the Bayes classifier perform so poorly?}
This initially made me curious about why the Bayes classifier performed so poorly, and compared Mahalanobis distance with the case III classifier; This
comparison is shown in Table \ref{tab:mahalanobis_caseIII}. The case III classifier, which assumes each Gaussian has different shapes and sizes, performed significantly better
than the Mahalanobis distance approach. I suspect that although PCA does de-correlate data across both classes, the within-class covariance matrix may not be well
approximated with a diagonal matrix with the off-diagonal elements set to zero, which is what the Mahalanobis distance approach assumes. This could explain the significant
performance gap between the two methods.

\paragraph{Val/Test Error Gap.}
One interesting observation from the results in Table \ref{tab:val_test_gap} is the significant gap between validation and test errors for some of the SVM configurations, particularly for the
polynomial kernel at low resolution. For example, in fold 2 of the 16x20 resolution, the polynomial kernel achieved a validation error of 4.51\% but a test error
of 9.02\%, which is a substantial increase. This suggests that the model may be overfitting to the validation set, and highlights the importance of using a separate
test set to evaluate the true generalization performance of the model. Similar patterns can be observed in fold 2 of the 48x60 resolution for both polynomial and RBF
kernels, where validation errors were around 3.76\% but test errors jumped to around 9\%.

\begin{table}[h]
	\centering
	\caption{SVM and Bayes Classifier Results for Resolution 16x20}
	\begin{tabular}{clcccc}
		\toprule
		\textbf{Fold} & \textbf{Method} & \textbf{Best Parameters} & \textbf{Val Error (\%)} & \textbf{Test Error (\%)} \\
		\midrule
		\multirow{3}{*}{1}
		              & Poly            & $C=0.1,\ D=1$            & 8.27                    & 9.02                     \\
		              & RBF             & $C=1,\ \gamma=0.1$       & 9.02                    & 6.77                     \\
		              & Bayes           & ---                      & ---                     & 21.80                    \\
		\midrule
		\multirow{3}{*}{2}
		              & Poly            & $C=0.1,\ D=1$            & 4.51                    & 9.02                     \\
		              & RBF             & $C=1,\ \gamma=0.1$       & 6.02                    & 8.27                     \\
		              & Bayes           & ---                      & ---                     & 14.29                    \\
		\midrule
		\multirow{3}{*}{3}
		              & Poly            & $C=0.1,\ D=1$            & 7.52                    & 10.53                    \\
		              & RBF             & $C=10,\ \gamma=0.1$      & 6.77                    & 8.27                     \\
		              & Bayes           & ---                      & ---                     & 18.05                    \\
		\midrule
		\multirow{3}{*}{Avg}
		              & Poly            & ---                      & ---                     & 9.52                     \\
		              & RBF             & ---                      & ---                     & 7.77                     \\
		              & Bayes           & ---                      & ---                     & 18.05                    \\
		\bottomrule
	\end{tabular}
	\label{tab:svm_bayes_16x20}
\end{table}

\begin{table}[h]
	\centering
	\caption{SVM and Bayes Classifier Results for Resolution 48x60}
	\begin{tabular}{clcccc}
		\toprule
		\textbf{Fold} & \textbf{Method} & \textbf{Best Parameters} & \textbf{Val Error (\%)} & \textbf{Test Error (\%)} \\
		\midrule
		\multirow{3}{*}{1}
		              & Poly            & $C=0.1,\ D=3$            & 6.77                    & 6.77                     \\
		              & RBF             & $C=10,\ \gamma=0.1$      & 9.02                    & 8.27                     \\
		              & Bayes           & ---                      & ---                     & 18.80                    \\
		\midrule
		\multirow{3}{*}{2}
		              & Poly            & $C=0.1,\ D=1$            & 3.76                    & 9.02                     \\
		              & RBF             & $C=1,\ \gamma=0.1$       & 3.76                    & 8.27                     \\
		              & Bayes           & ---                      & ---                     & 11.28                    \\
		\midrule
		\multirow{3}{*}{3}
		              & Poly            & $C=0.1,\ D=1$            & 10.53                   & 9.77                     \\
		              & RBF             & $C=1,\ \gamma=0.1$       & 9.02                    & 10.53                    \\
		              & Bayes           & ---                      & ---                     & 12.03                    \\
		\midrule
		\multirow{3}{*}{Avg}
		              & Poly            & ---                      & ---                     & 8.52                     \\
		              & RBF             & ---                      & ---                     & 9.02                     \\
		              & Bayes           & ---                      & ---                     & 14.04                    \\
		\bottomrule
	\end{tabular}
	\label{tab:svm_bayes_48x60}
\end{table}

\begin{table}[h]
	\centering
	\caption{Selected SVM results highlighting val/test error gap}
	\begin{tabular}{llcccc}
		\toprule
		\textbf{Resolution} & \textbf{Method} & \textbf{Fold} & \textbf{Best Parameters} & \textbf{Val Error (\%)} & \textbf{Test Error (\%)} \\
		\midrule
		16x20               & Poly            & 2             & $C=0.1,\ D=1$            & 4.51                    & 9.02                     \\
		48x60               & Poly            & 2             & $C=0.1,\ D=1$            & 3.76                    & 9.02                     \\
		48x60               & RBF             & 2             & $C=1,\ \gamma=0.1$       & 3.76                    & 8.27                     \\
		\bottomrule
	\end{tabular}
	\label{tab:val_test_gap}
\end{table}

\begin{table}[h]
	\centering
	\caption{Mahalanobis Distance vs Case III. Best results are bolded, and second-best results are underlined.}
	\begin{tabular}{lcccc}
		\toprule
		              & \multicolumn{2}{c}{\textbf{16x20}} & \multicolumn{2}{c}{\textbf{48x60}}                                                      \\
		\cmidrule(lr){2-3} \cmidrule(lr){4-5}
		\textbf{Fold} & \textbf{Mahalanobis (\%)}          & \textbf{Case III (\%)}             & \textbf{Mahalanobis (\%)} & \textbf{Case III (\%)} \\
		\midrule
		1             & 21.8                               & \underline{13.53}                  & 18.80                     & \textbf{12.03}         \\
		2             & 14.29                              & 13.53                              & \textbf{11.28}            & \textbf{11.28}         \\
		3             & 18.05                              & \underline{12.03}                  & \underline{12.03}         & \textbf{11.28}         \\
		\midrule
		Avg           & 18.05                              & \underline{13.03}                  & 14.04                     & \textbf{11.53}         \\
		\bottomrule
	\end{tabular}
	\label{tab:mahalanobis_caseIII}
\end{table}

\end{document}
